%% ------------------------------------------------------------
%%  قلم‌ها
%%  قلم فارسی: Estedad (مجوز آزاد OFL) — داخل پوشهٔ fonts/ قرار دارد
%%  تا پروژه روی هر سیستمی (و Overleaf) بدون نصب قلم کامپایل شود.
%%  برای تغییر قلم کافی است فقط همین فایل را ویرایش کنید؛ مثلاً:
%%     \settextfont{XB Niloofar}   یا   \settextfont{Vazirmatn}
%% ------------------------------------------------------------
\usepackage{xepersian}

\settextfont[
  Path=fonts/, Extension=.ttf,
  UprightFont=*-Regular, BoldFont=*-Bold,
  Scale=0.92
]{Estedad}

% قلم‌های کمکی فارسی (برای عنوان‌ها و متن‌های ظریف)
\defpersianfont\headfont[
  Path=fonts/, Extension=.ttf,
  UprightFont=*-ExtraBold, BoldFont=*-Black
]{Estedad}
\defpersianfont\semifont[
  Path=fonts/, Extension=.ttf,
  UprightFont=*-SemiBold, BoldFont=*-Bold, Scale=0.92
]{Estedad}
\defpersianfont\lightfont[
  Path=fonts/, Extension=.ttf,
  UprightFont=*-Light, BoldFont=*-Medium, Scale=0.92
]{Estedad}

% قلم لاتین (متن انگلیسی داخل \lr{...}) — Latin Modern همراه همهٔ توزیع‌ها هست
\setlatintextfont{Latin Modern Roman}

% جداکنندهٔ شماره‌ها (مثلاً «تعریف ۲.۳»)
\SepMark{.}

% فاصلهٔ بین خطوط برای خوانایی بهتر متن فارسی
\linespread{1.32}
